\documentclass[10pt,a4paper]{article}
\usepackage{luatexja}
\usepackage[haranoaji]{luatexja-preset}
\usepackage{geometry}
\geometry{top=14mm,bottom=14mm,left=18mm,right=18mm}
\usepackage{amsmath,amssymb}
\usepackage{booktabs}
\usepackage{enumitem}
\setlist{nosep,leftmargin=*}
\usepackage{titlesec}
\titlespacing*{\section}{0pt}{6pt}{3pt}
\titleformat{\section}{\normalfont\bfseries\large}{}{0pt}{}
\usepackage{xcolor}
\usepackage{hyperref}
\hypersetup{colorlinks=true,linkcolor=black,urlcolor=blue,citecolor=black}
\renewcommand{\tablename}{表}
\pagestyle{empty}
\setlength{\parindent}{0pt}
\setlength{\parskip}{3pt}

\begin{document}

{\LARGE\bfseries 進捗の要約（2026-10-07）}\hfill 下沢\,亮太郎

\medskip
\textbf{目標}: 要件を入れると最適化プログラムを書く，最適化専用の小型AIエージェント（2027年3月発表）．
柱はプロンプト最適化（GEPA）と，プロンプト最適化とfine-tuningを交互に回す方法（BetterTogether）．

\section{今回やったこと}
\begin{itemize}
  \item \textbf{難しい問題集を作った．} これまでの問題は厳密解が出る小規模なもので，AIがほぼ満点に達していた．新たに受け取った28問を元に，厳密解が出ない規模の「大規模問題」を20種別$\times$40問，計800問生成し，学習用600問とテスト用120問に分けた．
  \item \textbf{正解を確かめた教師データで小型モデルを学習した．} 大きなAI（Claude，Qwen）が書いたプログラムのうち，実際に動かして正解と確かめたものだけを16,228件集め，Gemma 4 12Bと4Bを学習した．
  \item \textbf{プロンプト最適化と学習を交互に回した．} 学習済みモデルから始める周と，素のモデルから始める周の2通り．
\end{itemize}

\section{結果}
テストは学習に使っていない問題．正解 = 制約違反0で参照解から+10\%以内．

\begin{table}[h]
  \centering
  \small
  \begin{tabular}{lrrl}
    \toprule
    モデル & 大規模テスト120問 & 大規模の元問題28問 & 学習で見ていない種別 \\
    \midrule
    素のGemma 4 12B & 0 & 3 & 11問中3 \\
    Qwen3.6 27B & 4 & --- & --- \\
    Gemma 4 12B（学習後） & \textbf{93} & \textbf{14} & 11問中2 \\
    4B（学習後，12Bの1/3の大きさ） & 81 & 10 & 11問中3 \\
    \bottomrule
  \end{tabular}
\end{table}

\section{わかったこと}
\begin{enumerate}
  \item \textbf{学習は効く．} 大規模問題を1問も解けなかったGemma 12Bが120問中93問を解き，4Bでも81問．12Bの9割強を1/3の大きさで出せる．
  \item \textbf{ただし「学習で見た種類」に限る．} 学習データのなかった3種別はテスト18問中0問で，学習していない種別には転移しない．当面は対象を学習した20種別に絞り，Gemma 4をその専用エンジンとして仕上げる．
  \item \textbf{プロンプト最適化は素のモデルにしか効かない．} 素のGemmaではGEPAが指示文を改善し，別の問題集で+19問（345$\rightarrow$364）．しかし学習済みモデルは学習時の指示文に張り付いていて，GEPAは52案を試して1案も採用できなかった．改善した指示文で学習し直しても，結果は変わらなかった．「プロンプト最適化と学習の相乗効果」は，この問題集ではまだ観測できていない．
  \item \textbf{教師候補のGLM-5.3-Flashは不採用．} 手元のGPUでは1問40分かかり，3問で正解0．
\end{enumerate}

\section{次にやること}
\begin{itemize}
  \item 「生成$\rightarrow$実行$\rightarrow$検証$\rightarrow$修復」の流れをモデルに組み込み，単発の生成ではなくエージェントとして測る．プロンプト最適化は修復の段に限って使う．
  \item 20種別すべての学習データでGemma 4を学習し直し，対応外の要件を判定して弾く経路を作る．
  \item 大きなAIを直接使う場合と比べた，1問あたりの時間と費用の表を作る．
  \item \textbf{お願い}: 3月の評価用に，対応20種別の範囲で発表まで非公開の問題を5〜10問ご用意いただけないか．評価の手順は事前に合意したい．
\end{itemize}

\end{document}
